% Screen Reader Implementation - Code Examples

This section presents the technical implementation of screen reader support through the Web Speech API and ARIA live regions, enabling auditory feedback for blind and visually impaired users. The implementation employs a custom \texttt{useScreenReader} hook that provides consistent announcement patterns across all interactive components, ensuring compliance with WCAG 2.2 Success Criterion 4.1.3 (Status Messages) \cite{WCAG22}.

\subsubsection{Interactive Components (Button, Link, Toggle)}

Interactive button, link, and toggle components require descriptive labels that convey their purpose to assistive technology users. The implementation utilizes the \texttt{useScreenReader} hook to provide auditory feedback during user interactions.

\begin{lstlisting}[caption={Interactive components with screen reader support}, label={lst:interactive-sr}]
// Files: Button.jsx, Link.jsx, Toggle/index.js
const { speak } = useScreenReader();

<button onFocus={() => speak('Create new app button')} 
        aria-label="Create new application">Create App</button>

<Link to="/signup" onFocus={() => speak('Sign up link')} 
      aria-label="Create a new account">Sign up</Link>

<input type="checkbox" onFocus={() => speak(`${label}, ${checked ? 'enabled' : 'disabled'}`)} 
       aria-label={label} />
\end{lstlisting}

The \texttt{aria-label} attribute provides textual descriptions that assistive technologies announce automatically. The \texttt{onFocus} event handler invokes the \texttt{speak()} function to deliver auditory feedback, with toggles announcing both label and current state (e.g., "Dark mode, enabled").

\subsubsection{Text Input Fields}

Text input fields facilitate information collection from users. Assistive technologies announce the field's purpose and any existing values to provide context regarding the expected input data.

\begin{lstlisting}[caption={Text input field with value announcement}, label={lst:input-sr}]
// File: frontend/src/_ui/Input/index.js
const { speak } = useScreenReader();

const handleFocus = () => {
  const announcement = value 
    ? `${label}, current value: ${value}` 
    : label;
  speak(announcement);
};

<input type="text" value={value} 
       onFocus={handleFocus} aria-label={label} />
\end{lstlisting}

Upon focusing an input field, the screen reader announces the field name and any existing value, providing contextual information prior to editing operations.

\subsubsection{Radio Button Groups}

Radio button groups facilitate single-option selection from multiple choices. Assistive technologies require announcements of both the group purpose and individual option labels for effective selection operations.

\begin{lstlisting}[caption={Radio button group with ARIA roles}, label={lst:radio-sr}]
// File: Radio.jsx
<div role="radiogroup" aria-label="Select account type">
  {options.map(option => (
    <label>
      <input type="radio" aria-label={option.name} />
      {option.name}
    </label>
  ))}
</div>
\end{lstlisting}

The \texttt{role="radiogroup"} attribute communicates to assistive technologies that these input elements constitute a related option set. Each radio button's \texttt{aria-label} provides its individual label, enabling users to comprehend and select options effectively.

\subsubsection{Keyboard Navigation Shortcuts}

Keyboard shortcuts facilitate navigation and interaction without pointer device input. Assistive technologies announce available shortcuts and provide feedback during navigation operations.

\begin{lstlisting}[caption={Keyboard shortcut announcements}, label={lst:shortcuts-sr}]
// File: KeyboardNavigation.jsx
const { speak } = useScreenReader();

useEffect(() => speak('Use Ctrl + arrow keys to switch between pages'), []);

const handleKeyDown = (e) => {
  if (e.key === 'ArrowDown') speak('Next element');
};
\end{lstlisting}

The implementation announces available keyboard shortcuts upon page initialization and provides continuous auditory feedback as users navigate through interface elements via arrow key input.

\subsubsection{Form Validation and Error Messages}

Form validation errors require clear auditory announcements to enable assistive technology users to identify and correct input mistakes. The system employs high-priority speech synthesis for immediate error notification.

\begin{lstlisting}[caption={Form validation error announcements}, label={lst:validation-sr}]
// File: LoginForm.jsx
const { speak, announceError } = useScreenReader();

const handleSubmit = (formData) => {
  if (!validateEmail(formData.email)) {
    announceError('Invalid email address. Please enter a valid email.');
  }
};
\end{lstlisting}

Validation errors utilize the \texttt{announceError()} function with assertive priority to immediately notify users of input problems, while \texttt{speak()} provides confirmatory feedback during form submission operations.

\subsubsection{Core Accessibility Hook}

The \texttt{useScreenReader} hook provides text-to-speech capabilities throughout the application, establishing consistent announcement patterns across all interface components.

\begin{lstlisting}[caption={Core useScreenReader hook implementation}, label={lst:hook-sr}]
// File: useScreenReader.js
const useScreenReader = () => {
  const speak = (text, priority = 'polite') => {
    const announcement = document.createElement('div');
    announcement.setAttribute('role', 'status');
    announcement.setAttribute('aria-live', priority);
    announcement.textContent = text;
    document.body.appendChild(announcement);
    setTimeout(() => announcement.remove(), 1000);
  };
  return { speak };
};
\end{lstlisting}

This hook integrates browser speech synthesis with ARIA live regions \cite{WCAG22}. The \texttt{priority} parameter controls announcement urgency: \texttt{'polite'} queues announcements sequentially, while \texttt{'assertive'} interrupts current announcements for critical messages.

\subsubsection{Focus Announcer System}

The focus announcer system provides specialized announcement functions for diverse interaction contexts, establishing consistent and contextual feedback throughout the application interface.

\begin{lstlisting}[caption={Focus announcement functions}, label={lst:focus-announcer}]
// File: useScreenReader.js
const announceFocus = (elementName, elementType) => {
  speak(`Focused on ${elementName} ${elementType}`);
};

const announceActivation = (elementName, action) => {
  speak(`${elementName} ${action || 'activated'}`);
};

return { speak, announceFocus, announceActivation };
\end{lstlisting}

The \texttt{announceFocus()} function provides standardized focus announcements across interface components, while \texttt{announceActivation()} communicates user actions such as button activation or menu selection.

\subsubsection{Component Manipulation Shortcuts}

Component manipulation operations including resizing, copying, and deletion require auditory feedback to inform assistive technology users about the outcomes of keyboard shortcut actions.

\begin{lstlisting}[caption={Component resize announcements}, label={lst:resize-sr}]
// File: WidgetWrapper.jsx
const { speak } = useScreenReader();

const handleKeyDown = (e) => {
  if (e.key === '=' && isArrowKey) {
    speak(`Expanding ${componentName} ${direction}`);
  }
};
\end{lstlisting}

Upon pressing the \texttt{=} key in conjunction with arrow keys, the component dimensions expand in the specified direction. The screen reader announces the operation, such as "Expanding button component down".

\begin{lstlisting}[caption={Component selection announcements}, label={lst:selection-sr}]
// File: WidgetWrapper.jsx
const handleKeyDown = (e) => {
  if (e.key === 'Enter') {
    speak(`${componentName} ${componentType} selected`);
  } else if (e.key === 'Escape') {
    speak('Component deselected');
  }
};
\end{lstlisting}

The Enter key initiates component selection and announces the operation, while Escape terminates selection and communicates the deselection.

\subsubsection{Tab Component Navigation}

Tab components organize content into switchable panels. Assistive technologies announce tab labels, selection states, and provide feedback when users transition between tabs.

\begin{lstlisting}[caption={Tab navigation with screen reader support}, label={lst:tabs-sr}]
// File: QueryManagerHeader.jsx
const { speak } = useScreenReader();

<p onClick={() => speak(`${tab.label} tab selected`)}
   onFocus={() => speak(`${tab.label} tab${activeTab === tab.id ? ', selected' : ''}`)}
   role="tab" aria-selected={activeTab === tab.id} />
\end{lstlisting}

Upon focusing a tab element, the screen reader announces its label and current selection state. Activation via Enter key announces the tab transition.

\subsubsection{Time Picker Component}

Time picker components facilitate temporal value selection. Assistive technologies announce the current time value upon focus and provide feedback when the time value changes.

\begin{lstlisting}[caption={Time picker with value announcements}, label={lst:timepicker-sr}]
// File: TimePicker.jsx
const { speak } = useScreenReader();

<ReactDatePicker
  onFocus={() => speak(`${paramLabel || 'Time'} picker, value: ${value || 'not set'}`)}
  onChange={(date) => speak(`Time set to ${moment(date).format('HH:mm')}`)}
/>
\end{lstlisting}

The time picker announces its current value upon receiving focus. When users select a new time value, the updated value is announced immediately.

\subsubsection{Code Editor Components}

Code editor components support single-line and multi-line code input with specialized assistive technology support for edit mode navigation and announcement.

\begin{lstlisting}[caption={Code editor navigation mode}, label={lst:codeeditor-sr}]
// File: SingleLineCodeEditor.jsx
const { speak } = useScreenReader();

const handleFocus = () => {
  speak(`${paramLabel || 'code'}, press Enter to edit, Escape to exit`);
};

const handleKeyDown = (e) => {
  if (e.key === 'Enter') speak(`Entering ${paramLabel}`);
  else if (e.key === 'Escape') speak(`Exiting ${paramLabel}`);
};
\end{lstlisting}

Code editors implement a navigation mode that prevents inadvertent modifications during keyboard navigation. Upon receiving focus, the editor announces instructions to press Enter for editing and Escape to exit edit mode.

\subsubsection{Number Input Component}

Number input components facilitate numeric value entry with assistive technology announcements for current values and focus states.

\begin{lstlisting}[caption={Number input with value announcement}, label={lst:number-sr}]
// File: NumberInput.jsx
const { speak } = useScreenReader();

<input type="number" value={value} 
       onFocus={() => speak(`${paramLabel || 'number'} input, value: ${value || 'empty'}`)} />
\end{lstlisting}

Number inputs announce their label and current value upon receiving focus, providing contextual information prior to editing operations.

\subsubsection{Slider Component}

Slider components enable value selection within a defined range via a draggable handle. Assistive technologies announce the current value as a percentage and provide feedback when the value changes.

\begin{lstlisting}[caption={Slider with percentage announcements}, label={lst:slider-sr}]
// File: Slider.jsx
const { speak } = useScreenReader();

<Slider.Root value={[sliderValue]}
  onFocus={() => speak(`${paramLabel || 'slider'}, value: ${sliderValue} percent`)}
  onValueCommit={(v) => speak(`${paramLabel} set to ${v} percent`)} />
\end{lstlisting}

The slider announces its current value as a percentage upon receiving focus. When users adjust the slider position, the new value is announced immediately upon release.

\subsubsection{Textarea Component}

Textarea components provide multi-line text input with assistive technology support for announcing current content and field purpose.

\begin{lstlisting}[caption={Textarea with content announcement}, label={lst:textarea-sr}]
// File: Textarea/index.js
const { speak } = useScreenReader();
const label = props['aria-label'] || props.placeholder || 'textarea';

<textarea onFocus={() => speak(value ? `${label}, value: ${value}` : label)} />
\end{lstlisting}

Textarea components announce their label and current content upon receiving focus. Empty fields announce only the label.

\subsubsection{Query Action Buttons}

Query management operations require clear announcements to notify users when queries are being executed or previewed, providing feedback for potentially time-intensive operations.

\begin{lstlisting}[caption={Query action button announcements}, label={lst:query-sr}]
// File: QueryManagerHeader.jsx
const { speak } = useScreenReader();

<ButtonComponent onClick={() => speak('Running query')}
                 onFocus={() => speak('Run query button')}>Run</ButtonComponent>

<ButtonComponent onClick={() => speak('Previewing query')}
                 onFocus={() => speak('Preview query button')}>Preview</ButtonComponent>
\end{lstlisting}

Query action buttons announce their purpose upon receiving focus and provide immediate feedback upon activation, notifying users that processing operations have commenced.

\subsubsection{Implementation Impact and Standards Compliance}

The screen reader implementations documented herein enable blind and visually impaired users to navigate the application exclusively via keyboard input, comprehend the purpose and state of all interactive elements, and receive immediate auditory feedback for their actions. These implementations demonstrate compliance with Web Content Accessibility Guidelines (WCAG) 2.2 Level AA requirements \cite{WCAG22}, specifically addressing Success Criteria 2.1.1 (Keyboard), 2.4.3 (Focus Order), 2.4.7 (Focus Visible), 4.1.2 (Name, Role, Value), and 4.1.3 (Status Messages). The consistent application of the \texttt{useScreenReader} hook throughout the codebase establishes standardized announcement patterns, ensuring maintainable and predictable assistive technology support across all interface components.
